% ===================== RESTAURAR ESTILO ANTES DEL FRONTMATTER =====================
\pagestyle{corporativo}

% ===================== FRONTMATTER CON MEMBRETE FORZADO =====================
\begin{frontmatter}

	\title{\projecttitle}

	\author[inst1]{DMV SAS\corref{cor1}}
	\ead{[correo@empresa.com]}

	\cortext[cor1]{Autor correspondiente}
	\address[inst1]{DMV SAS, Cali, Colombia}

	% ===================== ABSTRACT CON MEMBRETE GARANTIZADO =====================
	\begin{abstract}
		\thispagestyle{corporativo}
		\small  % Reducir tamaño para que quepa en una página
		El presente estudio técnico corresponde a la evaluación de validación y verificación solicitada por \textbf{Ingredion S.A.} a \textbf{DMV SAS}, respecto al diseño del sistema de calentamiento y el fondo toriesférico de reemplazo del tanque de almacenamiento de glucosa Tag 53A-90A-0056 en la planta de Cali, Colombia, documentado en los planos mecánicos REV2DMI17160102 (jun. 2026) y la información operativa suministrada por Ingredion S.A. El fondo de reemplazo incluye una chaqueta de media caña con un área de transferencia de calor de 14~m$^2$, un flujo de agua de servicio de 57,7~m$^3$/h a 75~\textdegree C, y se encuentra actualmente en fase de construcción. DMV SAS ejecutó la evaluación técnica mediante cálculo analítico independiente basado en propiedades termofísicas validadas de la glucosa Globe 1130, complementada con simulación numérica CFD en COMSOL Multiphysics. Los cálculos, utilizando correlaciones de Churchill-Chu para convección natural y Sieder-Tate para flujo interno, arrojan coeficientes globales entre 17,9 y 36,2~W/(m$^2\cdot$\textdegree C) en el rango de temperaturas de operación consideradas, con una resistencia térmica del lado del producto que concentra aproximadamente el 98~\% del total. Las simulaciones CFD validaron el modelo analítico con una concordancia superior al 95~\%. El análisis de transferencia de calor transitorio indica que el sistema mantiene la glucosa por encima del límite inferior de despacho de 57~\textdegree C durante el ciclo oficial de cinco descargas diarias, alcanzando una temperatura mínima de 58,56~\textdegree C, siempre que se operen las condiciones de diseño documentadas por Ingredion S.A.: agua de calentamiento a 75~\textdegree C, aislamiento térmico de 50,8~mm de lana mineral y entrada de glucosa entre 55~\textdegree C y 60~\textdegree C. El estudio identifica la convección natural del jarabe viscoso como el mecanismo limitante del proceso y cuantifica el impacto de las pérdidas térmicas al ambiente, concluyendo que el aislamiento térmico reduce las pérdidas en un 91,6~\% respecto al caso sin aislamiento.
	\end{abstract}

	\begin{keyword}
		glucosa \sep transferencia de calor \sep media ca\~na \sep tanque toriesf\'erico \sep convecci\'on natural \sep COMSOL \sep coeficiente global
	\end{keyword}

\end{frontmatter}

% ===================== RESTAURAR ESTILO DESPU\'ES DEL FRONTMATTER =====================
\pagestyle{corporativo}
% Forzar en la p\'agina actual tambi\'en
\thispagestyle{corporativo}
